\section{Forbidden paths and arithmetic savings}\label{sec:paths}

We now prove the graph estimate stated in
\Cref{prop:graph-estimate}.  Its main input is a bound for closed
walks after deleting a sparse periodic set of vertices.  In this
section we define that set and establish the arithmetic estimates
used to discard exceptional walks.  The parameters and kernels are
those of \Cref{sec:graph-setup}.  In the graph proof, $B$ tends to
infinity with $h,\tau,C_0,T$ fixed.

Deleting vertices that support short prohibited paths is part of the
divisibility-graph method of \cite[Section~4]{HR21}, developed for
composite labels in \cite[Sections~7 and~13.1]{Pilatte26}.
We use primitive specifications with an extra prime coordinate and prove
the resulting density and arithmetic bounds for the present graph.

Put
\begin{equation}\label{eq:trace-parameters}
 \ell=2\lfloor B\rfloor,\qquad
 L=\lfloor B^{1-\rho}\rfloor,\qquad
 t_0=\lceil B^{4\eps}\rceil.
\end{equation}
A \emph{closed line} is a list of signs and labels
$\boldsymbol d=((\varepsilon_i,d_i))_{i=1}^{\ell}$, where
$\varepsilon_i\in\{-1,1\}$, $d_i\in\mathcal D$, and the offsets
\[
 b_0=0,\qquad b_j=\sum_{i=1}^j\varepsilon_i h d_i
 \quad(1\le j\le\ell)
\]
satisfy $b_\ell=0$.  Repeated offsets are allowed.  Define its
normalized product by
\begin{equation}\label{eq:trace-kernel}
 K_{\boldsymbol d}(n)
   =\prod_{i=1}^{\ell}
     \frac{g_{d_i}(n+b_{i-1},n+b_i)}{W(n+b_{i-1})}.
\end{equation}
At each $B$, all residue expectations below are on the finite product
probability space of \Cref{sec:graph-setup}.

\subsection{The deleted vertices}

\begin{definition}[Path specifications]\label{def:primitive-specification}
A specification consists of a path with distinct integer offsets
$z_0=0,z_1,\ldots,z_m$, where $1\le m\le L$ and
\[
 z_i-z_{i-1}=\varepsilon_i h e_i,
 \qquad e_i\in\mathcal D,\quad \varepsilon_i\in\{-1,1\},
\]
an extra prime $p\in\mathcal S$ dividing none of the $e_i$, and an
index $i_0\in\{1,\ldots,m\}$ such that
$p\mid z_m-z_{i_0-1}$.  It \emph{qualifies at $x$} when
\[
 p\mid x,\qquad e_i\mid x+z_{i-1}\quad(1\le i\le m).
\]
Its prime support $\Pi$ is the union of $\{p\}$ and the prime factors
of all its edge labels.  Its \emph{cylinder} is the set of starting
integers $x$ satisfying these qualification conditions.  The cylinder
is either empty or fixes one residue at each prime in $\Pi$; in the
latter case its measure is $\prod_{q\in\Pi}q^{-1}$.

When testing a contiguous subpath as a specification, translate its
offsets to start at zero.  If it starts at $z_a$, test qualification
at $x+z_a$.  For the
reversed subpath, use the same convention with its new initial
vertex.  The extra prime may be any member of $\Pi$ absent from the
subpath's edge labels.  All these tests are constant as $x$ ranges
over a nonempty original cylinder, since they use only its fixed
prime coordinates.

A specification is \emph{primitive} if its cylinder is nonempty and
no shorter nonempty contiguous subpath, in either orientation, can
be made into a specification, with an extra prime from $\Pi$, that
qualifies at its own initial vertex for $x$ in that cylinder.  Set
\[
 Y=\{x\in\Z:\text{no primitive specification qualifies at }x\}.
\]
\end{definition}

The set $Y$ is periodic modulo $\prod_{p\in\mathcal S}p$.
For a path translated to start at $x$, a prime $q$ is \emph{active}
at its $i$th vertex when $q\mid x+z_i$.

\begin{lemma}[Descent to a primitive path]\label{lem:primitive-descent}
If a specification qualifies, a primitive specification using only
its prime support qualifies at one of its path vertices, along a
contiguous subpath in one of the two orientations.
\end{lemma}

\begin{proof}
Full edge divisibility survives reversal: $e\mid x$ and
$x'=x\pm he$ imply $e\mid x'$.  For a candidate subpath all residue
tests are constant on the original cylinder, since they involve only
its fixed prime coordinates.  If the current specification is not
primitive, take a shorter qualifying contiguous subpath in an allowed
orientation.  Its support is contained in the preceding support.
The positive length strictly decreases, so the procedure terminates.
\end{proof}

We shall prove the following estimate.  Its density assertion is
proved immediately; the closed-line bound is completed in
\Cref{sec:trace-count}.

\begin{theorem}[High trace]\label{thm:high-trace}
For the divisor family, cutoffs, and kernels of
\Cref{sec:graph-setup}, the periodic set $Y$ in
\Cref{def:primitive-specification} satisfies, as $B\to\infty$,
\begin{align}
 \Pbb(n\notin Y)&\le P_0^{-1+o(1)},\label{eq:deleted-density}\\
 \sum_{\boldsymbol d}
 \left|\E K_{\boldsymbol d}(n)
       \prod_{j=0}^{\ell}\1_Y(n+b_j)\right|
   &\le B^{-(1+c_*)\ell}.\label{eq:high-trace}
\end{align}
The sum is over all closed lines of length $\ell$.  The asymptotic
bounds are uniform over the admissible divisor families and cutoffs,
with $h,\tau,C_0,T$ fixed.
\end{theorem}

\begin{lemma}[Density of the deleted set]\label{lem:bad-density}
The set $Y$ satisfies \eqref{eq:deleted-density}.
\end{lemma}

\begin{proof}
It suffices to count all qualifying specifications.  For fixed edge
labels, signs, and suffix, its displacement $D=z_m-z_{i_0-1}$ is
nonzero because the path vertices are distinct, and
$|D|\le Lh\tau H$.  Sum the extra prime first, keeping these numerical
edge labels fixed.  Since $p>P_0$,
\[
 \sum_{\substack{p\in\mathcal S\,,\ p\mid D}}
       \frac1p
 \le \frac{\log(|D|+2)}{P_0\log 2}
 \ll \frac{\log(Lh\tau H+2)}{P_0}.
\]
The support weight for the remaining distinct edge primes is the
product of their reciprocals.  There are at most $LJ$ edge-prime
slots.  Choosing the length, signs, suffix, factor counts, and a
partition of these slots into equality classes costs
$\exp(O(LJ\log B))$.  Indeed a partition of at most $LJ$ slots has
at most $(LJ)^{LJ}$ descriptions.  Summing each remaining class over
its prime values costs at most $v_C+v_Z$, or $1+v_C+v_Z$ as a uniform
upper bound.  At this point we may drop all bin and distinctness
restrictions in those positive sums.  The resulting bound is
\[
 \Pbb(n\notin Y)
 \le P_0^{-1}
    \exp\bigl(O(LJ\log B)+O(LJ\log\log B)+O(\log B)\bigr).
\]
Since $LJ\log B=O(B^{1-\rho}\log^2B)=o(\log P_0)$, this proves
the assertion.
\end{proof}

The deletion also provides useful constraints on the labels of a
primitive path.  The next lemma identifies coefficients that remain
invertible even when a prime divides several edge labels.

\begin{lemma}[Prime intervals and the last constant block]
\label{lem:primitive-properties}
For sufficiently large $B$, let a primitive specification have edge
labels $e_1,\ldots,e_m$, extra prime $p$, and suffix ending at $m$.
Then the following statements hold.
\begin{enumerate}
\item For every edge prime $q$, the indices $i$ with $q\mid e_i$
form an interval of consecutive indices.
\item Let $b$ be the first index of the final maximal constant block
$e_b=\cdots=e_m$.  Then $b>1$, the signs on this block are constant,
and there exists $r\mid e_b$ with $r\nmid e_{b-1}$.  Every such $r$
occurs only in the constant tail.  In the specified suffix
displacement, its coefficient as a prime variable is nonzero modulo
$p$.  This coefficient assertion holds for any edge prime whose
occurrences are confined to that tail.
\item If $q$ occurs before $b$ and $r$ is as in the preceding part,
the coefficient of $q$ in $z_{b-1}-z_0$ is nonzero modulo $r$.
\end{enumerate}
Coefficients here are obtained by fixing all other distinct prime
values; squarefreeness makes each displacement linear in the prime
being varied.
\end{lemma}

\begin{proof}
Suppose successive occurrence blocks of $q$ are separated by edges
not containing $q$.  The intervening path begins and ends at active
vertices for $q$, since adjacent qualifying $q$-edges have divisible
endpoints.  It is a shorter simple $q$-free path with nonzero total
displacement divisible by $q$.  With extra prime $q$ and its whole
path as suffix, it contradicts primitivity.  This proves the first
part.

Within a constant-label block, an adjacent change of sign would
repeat a vertex, so all signs agree.  If $b=1$, the prescribed suffix
displacement is $\pm khe_m$ for some $1\le k\le L$.  The prime $p$
divides neither $e_m$ nor $h$, and $p>L$, so this is impossible.
Thus $b>1$.  The unequal squarefree labels $e_b,e_{b-1}$ lie in one
bin of ratio less than $2$.  If every prime of $e_b$ divided
$e_{b-1}$, their quotient would be an integer at least $2$, a
contradiction.  Choose $r\mid e_b$ absent from $e_{b-1}$.  By the
first part it is absent from the entire prefix.  Its coefficient in
the suffix is
\[
 \pm kh\frac{e_b}{r},\qquad 1\le k\le L.
\]
This is nonzero modulo $p$, because $p>h,L$ and $p$ divides none of
the edge labels.  The same reasoning applies to every prime confined
to the tail.

Finally, intersect the occurrence interval of $q$ with the prefix
$1,\ldots,b-1$.  The resulting nonempty interval is $i,\ldots,j$;
in particular, $j=b-1$ if $q$ also occurs in the tail.  The sum of
these prefix terms is $S=z_j-z_{i-1}\ne0$.  If $r\mid S$, reverse the path from
$z_{b-1}$ to $z_{i-1}$.  Its edges avoid $r$, it starts at a vertex
active for $r$, and its suffix from $z_j$ to $z_{i-1}$ has
displacement divisible by $r$.  Full divisibility survives reversal.
This shorter qualifying specification contradicts primitivity.
Since $q\ne r$ and $S$ is $q$ times the asserted coefficient, that
coefficient is invertible modulo $r$.
\end{proof}

\subsection{Absolute counts and triangular constraints}

We next record a deliberately coarse counting estimate.  It is used
only when an additional arithmetic saving makes an entire class
negligible.  The main contribution in \Cref{sec:trace-count} will
require a more economical enumeration.

For a line occurrence $(i,p)$, $p\mid d_i$, call it \emph{lit} if
$p\mid n+b_{i-1}$ and \emph{unlit} otherwise.  The condition is the
same at the arrival vertex.  We may attach a specification at a
line vertex $n+b_j$; qualification then refers to that starting
point.  A \emph{slot pattern} records the factor counts of all edge
labels, their signs, their bands, equality classes of prime slots,
and the lengths, attachments, extra-prime slots, and suffix indices
of attached specifications.  Distinct classes represent distinct
prime variables.  The pattern may also specify line occurrence
statuses and a tree on the visited vertices whose edges are recorded
line steps.

\begin{lemma}[Absolute reciprocal counting]\label{lem:crude-count}
Consider a line of length $\ell$ together with at most $2t_0$
attached qualifying specifications, each of length at most $L$.
Refine its slot pattern by recording the lit or unlit status of every
line occurrence.
There are $\exp(O(B\log^2B))$ slot patterns, including signs,
attachments, occurrence statuses, and recorded vertex identifications.
For each assigned pattern and its numerical prime values, let
$\mathfrak A$ be its attached specifications and let $E$ be any
residue event imposing the recorded occurrence statuses, possibly
with further restrictions.  Then
\begin{equation}\label{eq:absolute-pattern-majorant}
 \E\left[|K_{\boldsymbol d}(n)|\1_E(n)
     \prod_{\sigma\in\mathfrak A}
       \1_{\{\sigma\text{ qualifies at its assigned vertex}\}}\right]
 \le A^{\ell J}\prod_{p\text{ used}}\frac1p.
\end{equation}
If the recorded statuses include at least $t_0$ unlit occurrences
among center primes occurring at least twice on the line, this majorant has the additional
factor $P_0^{-t_0/2}$.  The total contribution of that class, summed
over all patterns and numerical assignments, is
$O(\exp(-B^{1+\eps/2}))$.
\end{lemma}

\begin{proof}
The number of slots is at most
\[
 \ell J+2t_0(LJ+1)=O(B\log B),
\]
since $4\eps<\rho$.  A partition of $R$ slots has at most $R^R$
descriptions.  All further indices have polynomially many choices
in $B$ per slot or per recorded vertex, giving the stated entropy.
Recording the destination of each step among at most $\ell+1$
vertices includes possible vertex identifications and any chosen
tree of recorded steps within the same bound.

Drop the cutoffs and $W^{-1}$, which lie in $[0,1]$.  The powers of
$A$ from core occurrences are at most $A^{\ell J}$.  A lit occurrence
of $p$, or qualification of any attached specification using $p$, fixes its
residue and supplies at most $1/p$.  In the absence of either, $p$
appears only on unlit line occurrences.  An unlit core occurrence
vanishes; each unlit center occurrence supplies $\theta/p$, so at
least one factor $1/p$ is still available.  The remaining conditions
may be discarded after taking these positive primewise bounds.
Independence of the prime residue coordinates proves
\eqref{eq:absolute-pattern-majorant}.

If a repeated center prime has a lit occurrence, every unlit
occurrence supplies an additional reciprocal beyond the residue
probability.  If all its $m\ge2$ line occurrences are unlit, they
supply $p^{-m}$, saving at least $p^{-(m-1)}$ beyond the displayed
majorant; here $m-1\ge m/2$.  Additional qualification conditions can
only improve the estimate.  At least $t_0$ repeated-unlit occurrences
therefore save $P_0^{-t_0/2}$.

For each unrestricted prime variable its reciprocal sum is at most
$1+v_C+v_Z=O(1+\log B)$.  Thus the total before the extra saving is
$\exp(O(B\log^2B))$, including $A^{\ell J}$.  The saving has
logarithm at most $-\frac12 B^{1+3\eps}$, which proves the last
assertion for sufficiently large $B$.
\end{proof}

The next lemma is the arithmetic alternative to having many unlit
occurrences.  Its ordering hypothesis is essential: all selected
prime values cannot be summed independently without examining their
dependencies.

\begin{lemma}[Triangular arithmetic saving]\label{lem:triangular-saving}
In the positive count of \Cref{lem:crude-count}, suppose every
configuration under consideration supplies at least $B^{2\eps}$
distinct prime variables $q_1,\ldots,q_k$ and tests with the following
properties.  Test $j$ uses only $q_j$, the earlier variables
$q_1,\ldots,q_{j-1}$, and nonselected variables, and is either
\begin{enumerate}
\item $q_j\mid D_j$, where $D_j\ne0$ is independent of $q_j$ and
$\log(|D_j|+2)=O(B\log B)$; or
\item $a_jq_j+c_j\equiv0\pmod{r_j}$, where $a_j,c_j,r_j$ are
independent of $q_j$, $r_j$ is an earlier or nonselected prime
variable, and $a_j\not\equiv0\pmod{r_j}$.
\end{enumerate}
Assume the choices of tests and their order have
$\exp(O(B\log^2B))$ descriptions.  Then their total absolute
contribution, also allowing the attached specifications in
\Cref{lem:crude-count}, is
\begin{equation}\label{eq:negligible-constraints}
 O\bigl(\exp(-B^{1+\eps/2})\bigr).
\end{equation}
The description bound applies in particular when each test uses a
bounded number of displacements along contiguous line or attached
paths, or along paths in a recorded tree of line steps, together with
one common specified integer offset of size $O(\ell H)$.
\end{lemma}

\begin{proof}
Fix the nonselected variables.  For the first type of test,
\[
 \sum_{\substack{q\in\mathcal S\,,\ q\mid D_j}}\frac1q
 \le \frac{\log(|D_j|+2)}{P_0\log 2}
 \ll \frac{B\log B}{P_0}.
\]
For the second type, invertibility leaves one residue class modulo
$r_j$.  Comparison with the integral of $1/x$ along that progression
gives, even when the sum is enlarged from primes to integers,
\[
 \sum_{\substack{P_0<n\le e^B\\ n\equiv a_0\pmod{r_j}}}\frac1n
 \le \frac1{P_0}+\frac{B-\log P_0}{r_j}
 \le \frac{1+B}{P_0}.
\]
The bound includes the zero residue class.

Sum the selected variables in the order $q_k,q_{k-1},\ldots,q_1$.
When summing $q_j$, every remaining earlier test is independent of
it, and its own test has one of the preceding uniform bounds.  The
conditions $D_j\ne0$ and $a_j\not\equiv0\pmod{r_j}$ are retained
through this elimination.  If either fails for fixed earlier data,
the admissible inner sum is zero.  Numerical distinctness may also
be retained; dropping it is unnecessary for the progression bound.
Backward induction therefore gives a factor
\[
 \left(\frac{O(B\log B)}{P_0}\right)^k.
\]
The other prime sums and all descriptions cost
$\exp(O(B\log^2B))$.  Since
$k\log P_0\ge B^{1+\eps}$, their product satisfies
\eqref{eq:negligible-constraints}.

We verify the final description convention.  There are
$O(B\log B)$ slots and polynomially many recorded positions or
vertices.  A bounded number of paths per test is specified by its
endpoints and type, together with a choice of prime variables and
order; choosing at most the number of slots many tests costs
$\exp(O(B\log^2B))$.  The common integer offset has
$O(\ell H+1)=\exp(O(B))$ possible values.  After fixing it we may
drop the equation that originally identifies it with a difference
of tree offsets, since we are taking an upper bound.  Each edge
product under arbitrary slot assignments is at most $e^{BJ}$.
Hence every indicated bounded sum of path displacements still has
logarithmic size $O(B\log B)$, even after the original divisor-bin
restrictions have been relaxed.  Nonzero and invertibility tests
are never relaxed.
\end{proof}

These estimates dispose of any class with sufficiently many
independent arithmetic restrictions.  We next use the tree formed
by first visits to find those restrictions and describe what remains.
